% Replace the existing Section 2, including its heading, with this file.
% Requires amsmath and amssymb. All pr: labels are local to this replacement.
\section{Problem formulation}
\label{pr:sec:problem}

\subsection{Three-dimensional mission and threat environment}
\label{pr:sec:environment}
The mission is defined in a three-dimensional Cartesian workspace
$\mathcal{W}\subset\mathbb{R}^{3}$. The UAV position is
$\mathbf{p}=[x,y,z]^{\mathsf T}$, where $x$ and $y$ are horizontal
coordinates and $z$ is altitude. The prescribed mission endpoints are
$\mathbf{p}_{s}$ and $\mathbf{p}_{g}$. A discrete search path is written as
\begin{equation}
 \Gamma_d=(\mathbf{p}_0,\ldots,\mathbf{p}_N),\qquad
 \mathbf{p}_0=\mathbf{p}_{s},\quad \mathbf{p}_N=\mathbf{p}_{g},
 \label{pr:eq:discrete}
\end{equation}
whereas the final time-parameterized trajectory is
$\mathbf{r}(t)\in\mathcal{W}$, $t\in[0,T]$. These representations are
distinguished because collision-free graph connectivity does not establish
continuous-trajectory feasibility. The search, geometric refinement and
speed-assignment procedures are described in Section~3.

Hard exclusion regions and traversable soft-threat regions are modeled
separately. Each hard zone is a vertical cylinder,
\begin{equation}
 \Omega_{F,j}=\left\{\mathbf{p}\in\mathcal{W}:
 (x-x_{F,j})^2+(y-y_{F,j})^2\le R_{F,j}^{2},\quad
 0\le z\le H_{F,j}\right\},
 \label{pr:eq:hardzone}
\end{equation}
where $R_{F,j}$ and $H_{F,j}$ are its horizontal radius and height. With
$\Omega_{\mathrm{hard}}=\bigcup_{j=1}^{N_F}\Omega_{F,j}$, the continuous
collision-avoidance requirement is
\begin{equation}
 \mathbf{r}(t)\notin\Omega_{\mathrm{hard}},\qquad t\in[0,T].
 \label{pr:eq:collision}
\end{equation}

For radar $i$ at $(x_{R,i},y_{R,i},0)$, define the slant range
\begin{equation}
 d_i(\mathbf{p})=\sqrt{(x-x_{R,i})^2+(y-y_{R,i})^2+z^2}.
 \label{pr:eq:slant}
\end{equation}
Its coverage is hemispherical within the nonnegative-altitude workspace.
Let $R_i$ be its outer radius and $R_{\mathrm{in},i}=\eta R_i$, with
$0<\eta<1$. For the unit radar scaling factors used in the experiments,
the local exposure field is
\begin{equation}
 P_{d,i}(\mathbf{p})=
 \begin{cases}
  1, & d_i(\mathbf{p})\le R_{\mathrm{in},i},\\
  \left(\dfrac{R_i-d_i(\mathbf{p})}{R_i-R_{\mathrm{in},i}}\right)^{\nu},
       & R_{\mathrm{in},i}<d_i(\mathbf{p})<R_i,\\
  0, & d_i(\mathbf{p})\ge R_i,
 \end{cases}
 \label{pr:eq:exposurefield}
\end{equation}
where $\nu>0$ controls the decay rate. The experiments use $\eta=0.3$
and $\nu=4$. Overlapping fields are combined by
\begin{equation}
 P_d(\mathbf{p})=\max_{1\le i\le N_R}P_{d,i}(\mathbf{p}),
 \qquad 0\le P_d(\mathbf{p})\le1.
 \label{pr:eq:maxfield}
\end{equation}
This is a prescribed dominant-exposure surrogate, rather than a statistical
fusion formula for the probability of detection by at least one radar.
It does not model radar scan timing, track accumulation or measurement
dependence. Hard zones determine admissibility, while the soft field
contributes to the accumulated exposure objective in
Section~\ref{pr:sec:metrics}.

\subsection{Fixed-wing flight constraints}
\label{pr:sec:constraints}
The model assumes still air and constant aircraft mass. Hence the scalar
path speed $v(t)=\|\dot{\mathbf{r}}(t)\|$ is identified with airspeed.
Feasibility refers to the specified geometric and reduced flight-envelope
constraints; it does not imply validation of a six-degree-of-freedom
aircraft model or a closed-loop tracking controller.

\subsubsection{Mission altitude envelope}
Let $d_s(\mathbf{p})$ and $d_g(\mathbf{p})$ denote the horizontal distances
to the start and goal. For transition distance $D_{\mathrm{tr}}$, define
\begin{equation}
 \alpha_e(\mathbf{p})=\operatorname{clip}
 \left(\frac{d_e(\mathbf{p})}{D_{\mathrm{tr}}},0,1\right),
 \qquad e\in\{s,g\},
 \label{pr:eq:transition}
\end{equation}
and $h_e(\mathbf{p})=z_e+(h_{\mathrm{en}}-z_e)\alpha_e(\mathbf{p})$.
Here $h_{\mathrm{en}}$ is the minimum en-route altitude. The local envelope is
\begin{equation}
 h_{\min}(\mathbf{p})=\min\{h_s(\mathbf{p}),h_g(\mathbf{p}),h_{\mathrm{en}}\},
 \qquad h_{\min}(\mathbf{r}(t))\le z(t)\le h_{\max}.
 \label{pr:eq:altitude}
\end{equation}
This construction retains the endpoint altitudes while imposing the
en-route lower bound outside both transition regions. A sufficient
compatibility condition for the ramp slope is
\begin{equation}
 D_{\mathrm{tr}}\ge
 \max_{e\in\{s,g\}}\frac{\max(h_{\mathrm{en}}-z_e,0)}{\tan\gamma_{\max}}.
 \label{pr:eq:ramp}
\end{equation}
This condition concerns the altitude envelope itself; it does not guarantee
that a feasible route exists in the presence of obstacles. In the present
mission, $z_s=z_g=0$ is an altitude reference for airborne trajectory
segments. Ground roll, take-off dynamics and landing dynamics are not modeled.

\subsubsection{Airspeed constraint}
The aircraft must remain above its stall-speed lower bound. The speed
assignment uses a prescribed cruise ceiling $v_{\mathrm{cr}}$, additionally
bounded by the aircraft maximum speed $v_{\max}$:
\begin{equation}
 v_{\mathrm{stall}}\le v(t)\le v_{\mathrm{cap}},\qquad
 v_{\mathrm{cap}}=\min\{v_{\mathrm{cr}},v_{\max}\}.
 \label{pr:eq:speedcap}
\end{equation}
For the reported configuration, $v_{\mathrm{stall}}=33~\mathrm{m\,s^{-1}}$,
$v_{\mathrm{cr}}=55~\mathrm{m\,s^{-1}}$ and
$v_{\max}=83~\mathrm{m\,s^{-1}}$. The airborne initial condition is
$v(0)=v_{\mathrm{stall}}$. No terminal-speed equality is prescribed;
$v(T)$ follows from the local ceilings and the forward--backward speed
assignment. The current speed assignment seeks the available speed under
these ceilings and does not solve a separate fuel-optimal speed-control problem.

\subsubsection{Turning constraint}
For a coordinated level turn, the bank-limited radius satisfies
\begin{equation}
 R_{\min}(v)=\frac{v^2}{g\tan\phi_{\max}}.
 \label{pr:eq:turnradius}
\end{equation}
The implementation applies this relationship as a curvature-based envelope
test on the three-dimensional path. For a regular spline
$\mathbf{q}(u)$, its spatial curvature is
\begin{equation}
 \kappa(u)=\frac{\|\mathbf{q}'(u)\times\mathbf{q}''(u)\|}
                    {\|\mathbf{q}'(u)\|^3},\qquad
 v_{\mathrm{turn}}(u)=\sqrt{\frac{g\tan\phi_{\max}}{\kappa(u)}}.
 \label{pr:eq:curvature}
\end{equation}
The turning ceiling is interpreted as infinite where $\kappa=0$.
The assigned speed must not exceed both $v_{\mathrm{cap}}$ and
$v_{\mathrm{turn}}$. If $v_{\mathrm{turn}}<v_{\mathrm{stall}}$, the geometry
is rejected; the speed is not reduced below stall to make it acceptable.
Using spatial curvature in this level-turn relationship is a reduced-model
approximation for gently climbing trajectories, not an exact bank-angle
reconstruction for general three-dimensional flight. A constant stall
threshold is used; bank-dependent stall speed and full load-factor dynamics
are outside the present model.

\subsubsection{Climb and vertical-rate constraints}
With horizontal arc length $s_h$, the flight-path angle and vertical rate obey
\begin{equation}
 \gamma=\arctan\left(\frac{\mathrm{d}z}{\mathrm{d}s_h}\right),\qquad
 |\gamma|\le\gamma_{\max},\qquad
 -v_{\mathrm{des},\max}\le\dot z\le v_{\mathrm{asc},\max}.
 \label{pr:eq:vertical}
\end{equation}
The angle bound restricts geometry, while $\dot z=v\sin\gamma$ couples
that geometry to speed. Both are required. The current limits are
$\gamma_{\max}=15^\circ$, $v_{\mathrm{asc},\max}=12~\mathrm{m\,s^{-1}}$
and $v_{\mathrm{des},\max}=15~\mathrm{m\,s^{-1}}$.

\subsubsection{Longitudinal acceleration and numerical feasibility}
The tangential acceleration is bounded by
\begin{equation}
 -a_{\mathrm{dec},\max}\le\dot v(t)\le a_{\mathrm{acc},\max},
 \label{pr:eq:acceleration}
\end{equation}
with $a_{\mathrm{acc},\max}=3~\mathrm{m\,s^{-2}}$ and
$a_{\mathrm{dec},\max}=4~\mathrm{m\,s^{-2}}$ in the experiments.
These bounds are imposed during speed assignment rather than added to the cost.

Search edges are screened for environmental and geometric feasibility.
After spline construction, endpoint retention, terminal horizontal heading,
local deviation, altitude, obstacle clearance and curvature are checked
again, followed by speed-envelope validation. The terminal reference heading
is inferred from the incoming discrete path, rather than supplied as an
independent mission heading. The fine-search corridor restricts graph search;
final spline locality is checked against the input path, and membership in
the coarse corridor is diagnostic by default.

The continuous constraints above specify the intended model. Their
implementation uses finite-resolution sampling and numerical tolerances:
the spline geometry is probed at 4096 parameter values, endpoints are accepted
within $2~\mathrm{m}$, and the terminal-heading tolerance is $12^\circ$.
Therefore, acceptance denotes satisfaction of these numerical tests, not
a formal continuous-domain safety certificate. Numerical settings and
resolution checks must accompany the reported experiments.

\subsection{Time, threat-exposure, and fuel models}
\label{pr:sec:metrics}
The mission metrics are flight time $T$, accumulated exposure $C$, and
fuel consumption $F$. They are evaluated on a common trajectory and with
the same time increments. Let $s\in[0,L]$ denote three-dimensional arc
length along the final curve and let $v(s)>0$ be its assigned path speed.
The associated time law is $\mathrm{d}t=\mathrm{d}s/v(s)$.

\subsubsection{Flight time}
The continuous flight-time functional is
\begin{equation}
 T=\int_0^L\frac{1}{v(s)}\,\mathrm{d}s.
 \label{pr:eq:time}
\end{equation}
For final trajectory samples $\mathbf{q}_i$ with speeds $v_i$, the numerical
evaluator uses the chord length
$\ell_i=\|\mathbf{q}_{i+1}-\mathbf{q}_i\|$ and constant tangential
acceleration on each sampled segment. Consequently,
\begin{equation}
 \Delta t_i=\frac{2\ell_i}{v_i+v_{i+1}},\qquad
 a_i=\frac{v_{i+1}^2-v_i^2}{2\ell_i},\qquad
 T_{\mathrm{num}}=\sum_i\Delta t_i,
 \label{pr:eq:segmenttime}
\end{equation}
for nonzero-length segments. This is exact for the adopted segment motion
model, while the sampled representation approximates the continuous spline.
It avoids replacing the final trajectory speed by one mission-wide constant.

During graph search, the speed profile is not yet available. A straight
edge $e$ of length $\ell_e$ and altitude change $\Delta z_e$ instead uses
\begin{equation}
 \Delta t_e^{\mathrm{search}}=
 \max\left\{\frac{\ell_e}{v_{\mathrm{cr}}},
             \frac{|\Delta z_e|}{v_{z,\max}(e)}\right\},
 \label{pr:eq:searchtime}
\end{equation}
where $v_{z,\max}(e)$ is the ascent or descent limit according to the sign
of $\Delta z_e$; the second term is zero for level flight. Here
$v_{\mathrm{cr}}$ lies within the aircraft speed interval. This edge model
is a search surrogate and does not account for acceleration across graph
vertices. The final metrics are recomputed after geometric refinement and
speed assignment; they need not equal the search-stage estimates.

\subsubsection{Accumulated threat exposure}
The exposure objective is
\begin{equation}
 C=\int_0^T P_d(\mathbf{r}(t))\,\mathrm{d}t
   =\int_0^L\frac{P_d(\mathbf{q}(s))}{v(s)}\,\mathrm{d}s.
 \label{pr:eq:exposure}
\end{equation}
Because $P_d$ is dimensionless and bounded by unity, $C$ has units of
seconds and satisfies $0\le C\le T$. It measures time-weighted exposure,
not mission detection probability, interception probability or expected
loss. The same geometric passage can incur different exposure when the
time spent within the field changes. Thus a length integral of $P_d$ alone
would not represent the objective evaluated here.

\subsubsection{Fuel consumption}
Let $v_h=\sqrt{\dot x^2+\dot y^2}$ and $v_z=\dot z$ be the horizontal and
vertical velocity components. The implemented fuel model uses a
quasi-steady parabolic-drag power surrogate. Air density is approximated by
\begin{equation}
 \rho(z)=\rho_0\exp(-z/H_\rho),\qquad
 \rho_0=1.225~\mathrm{kg\,m^{-3}},\quad H_\rho=8500~\mathrm{m}.
 \label{pr:eq:density}
\end{equation}
This is an exponential atmosphere approximation. It is not the full
piecewise International Standard Atmosphere model. The level-flight
mechanical power surrogate is
\begin{equation}
 P_{\mathrm{lev}}(v_h,z)=
 \frac{1}{2}\rho(z)C_{D0}S v_h^3+
 \frac{2(mg)^2}{\rho(z)v_h S\pi\,\mathrm{AR}\,e_O},
 \label{pr:eq:levelpower}
\end{equation}
where $S$ is wing area, $C_{D0}$ is the zero-lift drag coefficient,
$\mathrm{AR}$ is aspect ratio, and $e_O$ is the Oswald efficiency factor.
The first and second terms represent parasitic and induced power,
respectively. The model uses horizontal speed in these terms; it does not
substitute path speed $v$ for $v_h$ in climbing segments.

For the current configuration, descent energy recovery is zero. With
climb efficiency $\eta_c$ and fuel-to-propulsive efficiency $\eta_p$,
the fuel-equivalent power is
\begin{equation}
 P_f(v_h,v_z,z)=\frac{P_{\mathrm{lev}}(v_h,z)+P_c(v_z)}{\eta_p},
 \qquad
 P_c(v_z)=
 \begin{cases}
  mgv_z/\eta_c, & v_z>\delta_z,\\
  0, & v_z\le\delta_z,
 \end{cases}
 \label{pr:eq:fuelpower}
\end{equation}
where $\delta_z=0.01~\mathrm{m\,s^{-1}}$ is the small vertical-rate
deadband used in evaluation. The routine numerically floors $v_h$ at
$v_{\mathrm{stall}}/2$; this floor is inactive for trajectories satisfying
the stated stall and flight-path-angle bounds. With $H_f$ denoting the fuel
lower heating value,
\begin{equation}
 F=\frac{1}{H_f}\int_0^T
 P_f(v_h(t),v_z(t),z(t))\,\mathrm{d}t,
 \qquad H_f=43\times10^6~\mathrm{J\,kg^{-1}}.
 \label{pr:eq:fuel}
\end{equation}
Thus $F$ is reported in kilograms. Since $P_f$ already includes division
by $\eta_p$, the efficiency must not be applied a second time when
converting energy to fuel mass.

For numerical integration, each sampled segment is subdivided at spatial
fractions $0=\lambda_0<\cdots<\lambda_K=1$. Constant tangential acceleration
gives
\begin{equation}
 v_{ij}=\sqrt{(1-\lambda_j)v_i^2+\lambda_j v_{i+1}^2},\qquad
 \Delta t_{ij}=\frac{2\ell_i(\lambda_{j+1}-\lambda_j)}{v_{ij}+v_{i,j+1}}.
 \label{pr:eq:subtime}
\end{equation}
At subsegment midpoint $\mathbf{q}_{ij}^{m}$, the horizontal and vertical
velocities are calculated from the corresponding displacement divided by
$\Delta t_{ij}$. The three numerical metrics are
\begin{equation}
 \begin{aligned}
 T_{\mathrm{num}}&=\sum_{i,j}\Delta t_{ij},\\
 C_{\mathrm{num}}&=\sum_{i,j}P_d(\mathbf{q}_{ij}^{m})\Delta t_{ij},\\
 F_{\mathrm{num}}&=\frac{1}{H_f}\sum_{i,j}
 P_f(v_{h,ij},v_{z,ij},z_{ij}^{m})\Delta t_{ij}.
 \end{aligned}
 \label{pr:eq:quadrature}
\end{equation}
Spatial midpoint evaluation of the nonlinear field and power terms is a
quadrature approximation. Search edges use the same field and power
functions, with their edge time distributed uniformly over the subsegments.

This fuel model is intended for internally consistent comparisons. It
neglects fuel-induced mass change, bank-dependent induced-drag corrections,
transient engine behavior and an explicit tangential-acceleration power
term. These assumptions limit claims about absolute fuel prediction.

\subsection{Anchor-based normalized mission-oriented objective}
\label{pr:sec:objective}
Time, exposure and fuel have different units and numerical scales. A fixed
reference is therefore constructed from three single-objective searches
on the same coarse graph, with identical mission endpoints, environment,
constraints and search-stage evaluation rules. Denote the resulting paths
by $\Gamma_T$, $\Gamma_C$ and $\Gamma_F$, corresponding to the weight
vectors $(1,0,0)$, $(0,1,0)$ and $(0,0,1)$ in $(T,C,F)$ order. These
are discrete search anchors; they are not independently optimized final
continuous trajectories.

Evaluating all three metrics on each anchor gives the payoff table
\begin{equation}
 A=\begin{bmatrix}
 T^{\mathrm{search}}(\Gamma_T)&C^{\mathrm{search}}(\Gamma_T)&F^{\mathrm{search}}(\Gamma_T)\\
 T^{\mathrm{search}}(\Gamma_C)&C^{\mathrm{search}}(\Gamma_C)&F^{\mathrm{search}}(\Gamma_C)\\
 T^{\mathrm{search}}(\Gamma_F)&C^{\mathrm{search}}(\Gamma_F)&F^{\mathrm{search}}(\Gamma_F)
 \end{bmatrix}.
 \label{pr:eq:payoff}
\end{equation}
For objective $k\in\{T,C,F\}$, let the lower reference be the corresponding
anchor's diagonal entry, and let the upper reference be the column maximum:
\begin{equation}
 b_k=A_{kk},\qquad n_k=\max_{a\in\{T,C,F\}}A_{ak},\qquad
 \Delta_k=n_k-b_k>0.
 \label{pr:eq:anchorscale}
\end{equation}
The upper reference is a payoff-table nadir approximation, not the
maximum over all feasible paths or an exact continuous Pareto nadir.
Degenerate reference ranges must be rejected before normalized planning;
they are not silently replaced by an arbitrary denominator.

For a metric vector $\mathbf{f}=(T,C,F)$, define
\begin{equation}
 \widehat T=\frac{T-b_T}{\Delta_T},\qquad
 \widehat C=\frac{C-b_C}{\Delta_C},\qquad
 \widehat F=\frac{F-b_F}{\Delta_F}.
 \label{pr:eq:normalize}
\end{equation}
The same reference values are retained across weights, coarse and fine
searches, and final-trajectory reporting within a fixed experiment.
The normalized values are not clipped to $[0,1]$: a refined or retimed
trajectory can fall outside the three-anchor reference ranges. This does
not invalidate the normalization, but it precludes interpreting the
reference values as exact bounds on final continuous performance.

Mission preference is expressed through
\begin{equation}
 J=w_T\widehat T+w_C\widehat C+w_F\widehat F,\qquad
 w_T,w_C,w_F\ge0,\quad w_T+w_C+w_F>0.
 \label{pr:eq:objective}
\end{equation}
The coefficients specify relative preferences and need not sum to one.
Multiplication of all weights by the same positive constant preserves
path ordering for a fixed model. During graph search, the additive
contribution of edge $e$ is
\begin{equation}
 c_e=\frac{w_T\Delta T_e}{\Delta_T}
     +\frac{w_C\Delta C_e}{\Delta_C}
     +\frac{w_F\Delta F_e}{\Delta_F}.
 \label{pr:eq:edgeobjective}
\end{equation}
The lower-reference offsets are applied once to the complete path:
\begin{equation}
 J^{\mathrm{search}}(\Gamma_d)=\sum_{e\in\Gamma_d}c_e
 -\left(\frac{w_Tb_T}{\Delta_T}+\frac{w_Cb_C}{\Delta_C}
                         +\frac{w_Fb_F}{\Delta_F}\right).
 \label{pr:eq:offset}
\end{equation}
The bracketed term is path independent. Subtracting the lower references
on every edge would instead introduce a spurious dependence on the number
of edges and would not implement Eq.~\eqref{pr:eq:objective}.

This scalarization directs discrete route selection. Spline candidates
are subsequently ranked with a separate geometric refinement score, and
the final $(T,C,F,J)$ values are recomputed using the assigned speed
profile. The sequential procedure is therefore not a joint global
optimization of route geometry, speed and fuel, nor does it establish
global Pareto optimality of the final trajectories.

The sensitivity experiment fixes $w_C=1$ and evaluates
$w_T\in\{0,0.1,\ldots,1.1\}$ and
$w_F\in\{0,0.1,\ldots,0.5\}$, giving 72 configurations.
The coefficients are swept rather than tuned by an additional optimizer.
For tractability, this sweep uses coarse search followed by the same
types of spline refinement, feasibility checks and final evaluation; it
does not repeat the fine-corridor search for all 72 configurations.
Consequently, it characterizes preference sensitivity of that coarse-search
pipeline. Claims about the full two-stage pipeline require separate
controlled comparisons and representative-case experiments that include
the fine-corridor search.
